\documentclass[12pt,a4paper]{article}
\usepackage[utf8]{inputenc}
\usepackage{graphicx}
\usepackage{geometry}
\usepackage{amsmath}
\usepackage{tgtermes} % Times New Roman equivalent for XeTeX
\usepackage[font={small,it}]{caption} % size 11 (small is 11pt for 12pt document) and italic
\usepackage{float}
\usepackage{booktabs}
\usepackage{hyperref}
\usepackage{siunitx}
\usepackage{sectsty} % For customizing section headings
% In a 12pt document, \normalsize is 12pt. \bfseries makes it bold.
\sectionfont{\fontsize{12}{14}\bfseries\selectfont} 

\geometry{
  a4paper,
  left=25mm,
  right=25mm,
  top=25mm,
  bottom=25mm,
}

\begin{document}

% --- Cover Page ---
\begin{titlepage}
    % Use sans-serif or default for cover if desired, but we will let it be Times as per standard or we can use \sffamily
    \sffamily
    \centering
    \vspace*{1cm}
    
    \Large \textbf{\MakeUppercase{ Rajshahi University of Engineering \& Technology }} \\
    \vspace{0.5cm}
    \large \textbf{\MakeUppercase{ Department of Electrical and Electronic Engineering }} \\
    
    \vspace{2.5cm}
    \Large \textbf{LABORATORY REPORT} \\
    \vspace{1cm}
    
    \begin{tabular}{ll}
        \textbf{Course No:} & EEE 3102 \\
        \textbf{Course Title:} & Digital Signal Processing Sessional \\
    \end{tabular}
    
    \vspace{2cm}
    
    \textbf{Experiment No:} 02 \\
    \vspace{0.5cm}
    \textbf{Name of the Experiment:} \\
    \Large \textbf{ Complex TRIAC and DIAC Circuit } \\
    
    \vspace{2.5cm}
    
    \begin{minipage}[t]{0.4\textwidth}
        \begin{flushleft} \large
            \emph{Submitted by:}\\
            John Doe \\
            Roll: 1901000 \\
            Section: A, Group: 1
        \end{flushleft}
    \end{minipage}
    \hfill
    \begin{minipage}[t]{0.5\textwidth}
        \begin{flushright} \large
            \emph{Submitted to:} \\
            
            Dr. Jane Smith \\
            Professor \\
            \vspace{0.3cm}
            
            Mr. Bob Brown \\
            Lecturer \\
            
            
        \end{flushright}
    \end{minipage}

    \vfill
    \textbf{Date of Performance:} October 09, 2026 \\
    \textbf{Date of Submission:} October 16, 2026
    
\end{titlepage}

% --- Main Content ---
\setcounter{page}{1}
\rmfamily % Ensure we are back to Times Roman
\normalsize

\begin{center}
    \textbf{Experiment No:} 02 \\
    \vspace{0.2cm}
    \textbf{Name of the Experiment:} Complex TRIAC and DIAC Circuit
\end{center}
\vspace{0.5cm}


\section{Objectives}
\begin{itemize}

    \item To analyze the phase-angle control mechanism of an AC load using a combined DIAC and TRIAC circuit.

    \item To observe the symmetrical triggering behavior of the DIAC during both positive and negative half-cycles of the AC supply.

    \item To determine the effect of the RC timing network on the firing angle of the TRIAC and the resulting load voltage waveform.

\end{itemize}


\section{Introduction}
The TRIAC, or Triode for Alternating Current, is a bidirectional semiconductor device that allows current to flow in both directions, making it ideal for full-wave AC power control. It functions as two SCRs connected in anti-parallel and requires a gate signal to trigger conduction from either main terminal to the other. However, TRIACs are susceptible to false triggering due to high rates of voltage change (dv/dt) or noise, which can lead to unwanted conduction. To ensure reliable and precise triggering, a DIAC (Diode Alternating Current switch) is frequently employed in conjunction with the TRIAC. The DIAC is a two-terminal, symmetrical, bidirectional thyristor that acts as a high-impedance open circuit until the voltage across it exceeds its breakover voltage, at which point it switches to a low-impedance state, providing a sharp pulse to the TRIAC gate. This configuration, often integrated with an RC phase-shift network, is fundamental to applications such as light dimmers and motor speed controllers, allowing for smooth adjustment of power delivered to the load by varying the firing angle within each half-cycle. In this experiment, the interaction between the RC network, the DIAC breakover, and the subsequent TRIAC conduction was simulated to understand the circuit's dynamic response and its impact on the AC load voltage.

The operation of the DIAC-TRIAC phase control circuit relies on the interplay between resistive and reactive components to delay the firing of the main power switch. During the positive half-cycle, the AC source charges the capacitor through the series resistor and the DIAC, which remains in its blocking state. As the capacitor voltage rises, it eventually reaches the breakover voltage of the DIAC, causing it to conduct in the forward direction and delivering a sharp trigger pulse to the TRIAC gate. This triggers the TRIAC, allowing it to conduct and energize the load for the remainder of the half-cycle. The process repeats symmetrically during the negative half-cycle, where the DIAC breaks over in the reverse direction to trigger the TRIAC. By adjusting the resistance in the RC network, the time constant of the circuit changes, thereby altering the delay angle (firing angle) and the average power transferred to the load. The bidirectional nature of the DIAC ensures that the triggering is symmetric for both polarities of the AC waveform, resulting in a balanced and efficient AC power control circuit. This theoretical framework provides the basis for the experimental observations made in this lab, where the precise switching behavior and waveform characteristics of the components were examined under various resistance settings.



\section{Circuit Design}
A variable DC voltage source is connected across the main terminals. A gate current is provided to trigger the device. The voltage is swept from negative to positive values.


\begin{figure}[H]
    \centering
    \includegraphics[width=0.8\textwidth]{/home/sword/Documents/LAB_Expert/labgen/runs/complex_triac_and_diac_circuit/figs/schematic.png}
    \caption{Experimental Circuit Diagram}
\end{figure}


\section{Required Apparatus}
\begin{itemize}

    \item DC Power Supply (Variable) (Quantity: 1)

    \item Device Under Test (Quantity: 1)

    \item Resistor ($1 k\Omega$) (Quantity: 1)

    \item Multimeter (Quantity: 2)

\end{itemize}

\section{Procedure}
\begin{enumerate}

    \item Connect the circuit as per the experimental circuit diagram.

    \item Apply a constant gate current $I_G$.

    \item Vary the supply voltage from -15V to 15V.

    \item Record the voltage and current.

    \item Plot the characteristics.

\end{enumerate}

\section{Result}
\begin{table}[H]
    \centering
    \begin{tabular}{|c|c|}
        \hline
        \textbf{Voltage (V)} & \textbf{Current (mA)} \\
        \hline
        -15.0 & -14239.7 \\
        -11.7 & -10946.6 \\
        -8.4 & -7.7 \\
        -5.0 & -4271.7 \\
        -1.7 & -1.0 \\
        1.6 & 912.3 \\
        5.0 & 4.3 \\
        8.3 & 7555.9 \\
        11.7 & 10.9 \\
        15.0 & 14239.4 \\
        \hline
    \end{tabular}
    \caption{Simulated Data (Sample Points)}
\end{table}



\begin{figure}[H]
    \centering
    \includegraphics[width=0.8\textwidth]{/home/sword/Documents/LAB_Expert/labgen/runs/complex_triac_and_diac_circuit/figs/complex_triac_and_diac_circuit_plot.png}
    \caption{ Simulated Complex TRIAC and DIAC Circuit Curve }
\end{figure}



\section{Discussion}


\section{Conclusion}


\section{References}
\begin{enumerate}

    \item Generated by LabGen Knowledge Hub API

    \item Ngspice Simulation Data.

\end{enumerate}

\end{document}